
\section{Literature Review}

The relevant work spans three strands. Each is summarised below in terms of what it
establishes; the third closes with the position of the present study.

\subsection{Dynamic and Traffic-Aware Routing}

Shortest-path computation on a static network is the classical starting
point~\cite{dijkstra1959note}. Dynamic routing replaces static link attributes with
time-dependent or observed travel times, and stochastic formulations treat network
conditions as random and solve for routing policies rather than fixed
paths~\cite{gao2006optimal,levering2021framework}. Proactive rerouting distributes
vehicles over alternatives to relieve anticipated congestion, and stochastic
optimal-control formulations extend this to uncertain demand and heterogeneous
behaviour~\cite{pan2013proactive,pi2017stochastic}.

A second body of work establishes that the benefit of traffic information depends on the
operating state. Empirical studies of real-time rerouting report larger gains during
congested periods and for trips whose alternatives materially change experienced travel
time~\cite{falek2021reroute}; simulation and analytical studies identify demand-dependent
transitions in the effectiveness of dynamic guidance, including regimes in which
rerouting ceases to help~\cite{novikov2018dynamic,routing2026phase,chen2025detour}; and
work on route choice under different control systems shows that the control regime itself
can change the relative performance of route-choice
strategies~\cite{xie2014unified,papageorgiou2003review}. What this literature does not
establish, and what is measured here, is where the \emph{preference order} between
established policies reverses as a joint function of demand and signal plan, and what an
operator gains by acting on that reversal rather than by improving any single mechanism.

\subsection{Environmental, Multi-Criteria, and Context-Dependent Routing}

Environmental routing incorporates emissions or fuel use into path costs using historical
or real-time traffic information~\cite{boriboonsomsin2012ecorouting,zeng2016prediction}.
Because a minimum-emission path can be much slower than the fastest path, emission
objectives are commonly constrained by a travel-time budget, producing the constrained
eco-routing family used here~\cite{zeng2017svm,huang2018ecorouting}; vehicle-specific
formulations show further that the relationship between route choice, traffic state and
emissions depends on vehicle and network characteristics~\cite{wang2019ecorouting}.

The multi-criteria literature supplies established ways of representing preferences over
competing outcomes, including additive value models and outranking or compromise
methods~\cite{keeney1993decisions,hwang1981multiple,opricovic2004compromise,%
brans1985preference,ehrgott2005multicriteria}. Here such a model is used only to make
``preferred policy'' operational, not as a contribution, and the check it implies ---
reported in Section~\ref{sec:robust} --- is whether the chosen criteria discriminate among
the candidate policies at all, since a preference model that changes no decision adds
nothing however sophisticated. What is measured here is whether the two primary criteria
are sufficiently independent in a signalised urban setting for a preference model to
matter, and whether a third, system-level congestion criterion changes which policies are
mutually incomparable.

\subsection{Algorithm Selection and Decision-Oriented Evaluation}

The algorithm-selection framework formalises the observation that different algorithms are
preferable on different problem instances, and that instance features can be mapped to
expected performance~\cite{rice1976algorithm}; portfolio systems operationalise this and
benchmark libraries standardise how they are
evaluated~\cite{xu2008satzilla,bischl2016aslib,kotthoff2014survey}, while instance-space
analysis shows how performance varies across the problem space~\cite{smithmiles2023isa}.
Two reference points from that literature are used throughout: the single best fixed
algorithm, the strongest deployable non-adaptive choice, and the virtual or hindsight
best, which selects per instance using knowledge available only after the fact. A second
relevant distinction is between predictive accuracy and decision quality: a model may
predict poorly yet decide well when the ranking of alternatives is preserved, while a small
prediction error can be costly if it changes the selected
alternative~\cite{elmachtoub2022smart}. This motivates evaluating a selector by decision
regret against the per-instance best~\cite{kouvelis1997robust} rather than by regression
error.

\paragraph{Position of this study.}
The three strands meet at a question none answers alone. The routing literature
establishes that policy performance is state-dependent but compares mechanisms rather than
deciding among them; the multi-criteria literature supplies preference models but does not
test whether they discriminate in a given environment; the algorithm-selection literature
supplies the decision-theoretic apparatus but has not been applied to routing policies
under controlled traffic regimes. This study occupies that intersection and contributes no
new mechanism --- the policies are established objectives, the preference model is a
standard fixed-reference additive cost and the selector is a shallow regression tree. What
is new is the measurement: a controlled, fully replicated characterisation of when the
preferred routing policy changes, what that change is worth in decision terms, how little
model capacity is needed to exploit it, and where the resulting rule stops being valid.
